% Generated draft: requires mathematical and source review.
\documentclass[11pt]{article}
\usepackage[margin=25mm]{geometry}
\usepackage{fontspec}
\setmainfont{Latin Modern Roman}
\usepackage{amsmath,amssymb,mathtools,unicode-math}
\setmathfont{Latin Modern Math}
\usepackage{enumitem,needspace,xurl,hyperref}
\setlength{\parindent}{0pt}
\setlength{\parskip}{5pt}
\setlength{\emergencystretch}{3em}
\newcommand{\N}{\mathbb{N}}
\newcommand{\Z}{\mathbb{Z}}
\newcommand{\Q}{\mathbb{Q}}
\newcommand{\R}{\mathbb{R}}
\newcommand{\C}{\mathbb{C}}
\newcommand{\F}{\mathbb{F}}
\newcommand{\A}{\mathbb{A}}
\newcommand{\PP}{\mathbb{P}}
\newcommand{\E}{\mathbb{E}}
\newcommand{\eps}{\varepsilon}
\newcommand{\dd}{\,\mathrm{d}}
\begin{document}
\section*{UnsolvedMath ID: 11000151}
\textbf{Source problem code:} catalogue entry
\subsection{Definitions and mathematical statement}
Let \(A_5=B_6\) be the six-strand braid group, with standard generators \(a_1,\ldots,a_5\), relations \(a_i a_{i+1}a_i=a_{i+1}a_i a_{i+1}\), and \(a_i a_j=a_j a_i\) for \(|i-j|\geq2\). Put
\[
 c=a_1a_2a_3a_4,\qquad
 h=a_5a_4a_3a_2a_1^2a_2a_3a_4a_5,
 \qquad Q=B_6/\langle\!\langle c^5h^{-1}\rangle\!\rangle.
\]
The flattened \(a^2_1\) is read as \(a_1^2\), consistent with the displayed ten-letter \(h\) and the source's word lengths; the typography should be confirmed.

\emph{Question.} If a positive word \(w\) in \(a_1,\ldots,a_5\) represents
\[
 u=\Delta_4^4=c^{10}\in Q,
\]
must its letter length be \(40\), \(30\), or \(20\), and must its ordered factorization be Hurwitz equivalent to one of
\[
 \Delta_4^4=(a_1a_2a_3a_4)^{10},\qquad
 \Delta_5^2=(a_1a_2a_3a_4a_5)^6,\qquad
 h^2=(a_5a_4a_3a_2a_1^2a_2a_3a_4a_5)^2,
\]
respectively? Positive means no inverse letters, and powers contribute their expanded number of letters. In a standard group-factorization convention, an elementary Hurwitz move replaces \((x,y)\) by \((xyx^{-1},x)\), or its inverse. Such moves preserve the ordered product and number of factors, so the three candidate forms have their indicated distinct lengths. The source refers to [1] for its convention; whether global simultaneous conjugation or restrictions on intermediate factors are included must be checked rather than added.

The source reports an affirmative related result for a closed surface: Siebert and Tian [1] prove that a genus-two symplectic Lefschetz fibration with no reducible fibers and transitive monodromy is holomorphic. It describes the monodromy factorization as a product of conjugates of factors Hurwitz equivalent to \(\Delta_5^2\) or \(h^2\), and explicitly says \(\Delta_4^4\) lacks transitive monodromy. It then discusses each factor becoming a boundary twist for a surface with boundary, suggesting only one such factor for a single boundary twist. The fibration base, precise boundary model, and comparison maps are absent. The source stresses that closed-surface equivalence is weaker, with larger equivalence classes, so that result does not settle the boundary-surface question posed here.
\subsection{Short English statement}
In the specified quotient of the six-strand braid group, classify positive factorizations of \(\Delta_4^4\): must their lengths be exactly \(40,30,20\) and their Hurwitz classes the three displayed ones? The closed genus-two fibration result is related context, with a weaker equivalence that need not answer the boundary case.
\subsection{Sources}
\begin{itemize}
\item Local UnsolvedMath registry, permanent ID 11000151, source code catalogue entry, unnumbered Question, PDF page 133. No individual proposer is supplied.
\item Farb, \emph{Problems on Mapping Class Groups and Related Topics}. Supplied URL: \url{https://www.math.uchicago.edu/~farb/papers/mcgbook.pdf}; recorded access 2026-07-29; extraction pdf-font-structure. The volume contains 329 explicitly typeset headings across chapters, with repeated labels disambiguated by page.
\item Reference [1] is cited for Hurwitz equivalence and the recent-to-the-source Siebert--Tian result. No full bibliographic entry, fibration-base hypotheses, or exact comparison theorem is provided.
\item Supplied status open and OPEN-TRIAGE. The research summary says only ``Truncated.'' Current primary-literature status and redistribution terms for the publicly accessible author/publisher copy remain NEEDS\_REVIEW.
\end{itemize}
\end{document}
